\scanpage{68}
\deleted{$(p,q,r)$ doit être parmi les \deleted{\unsure{trois}} syst. énumérés

a) $p,2,2$ : $\frac1p+\frac12+\frac12>\bigl(1+(1-\frac1s)\bigr)$ i.e. $\frac1p+\frac1s>1$ impossible car $\frac1p,\frac1s\leq\frac12$

b) $(p,3,2)$ $(p=3,4,5)$, $s\geq p$ $\deleted{\unsure{\frac1p+\frac13+\frac12|\frac1p+\frac1s+\frac12}}$
\[
\frac1s+\frac1p+\underbrace{\frac13+\frac12}_{5/6}>2\qquad \frac1s+\frac1p>\frac76\quad\text{impossible car }\frac1s,\frac1p\leq\frac13
\]}

\[
\frac1p+\frac1q>2-\underbrace{\left(\frac1r+\frac1s\right)}_{\substack{\leq1\\\text{car }\frac1r\leq\frac12\\\frac1s\leq\frac12}}\geq1
\quad\text{donc }\frac1p+\frac1q\geq1\quad\text{impossible car }\frac1p,\frac1q\leq\frac12
\]
\uline{Impossible} !

Car si $\operatorname{Card} I\geq4$, on \unsure{aura|aurait}
\[
\frac1{d_1}+\frac1{d_2}+\frac1{d_3}+\frac1{d_4}\deleted{\unsure{+\sum_{5\leq i\leq\nu}\frac1{d_i}}}>2\deleted{\unsure{-2}}+\sum_{5\leq i\leq\unsure{\nu|v}}\left(1-\frac1{d_i}\right)
\]
et on vient de voir que c'est impossible !\margin{$\unsure{\nu'|v'|v\text{.}}$\par $2$}

\deleted{\uline{Étude du cas} \deleted{$\operatorname{Card} I=2$}\inserted{$(N,N)$} : En caractéristique $0$, c'est le cas \unsure{kummerien}
\[
G\simeq\mu_N(k)\simeq\mathbb Z/N\mathbb Z\qquad\ldots
\]
\deleted{\unsure{Plus généralement}} \unsure{sim|sin}, il en est ainsi en car. $p$ première à $N$.

Si $p\mid N$, on trouve une \inserted{$p$-}extension \inserted{(de degré $p^r$, \unsure{tott|tolt} ramifiée et \unsure{isomorphe|séparable})} d'une $N'$-extension \unsure{kummerienne} $(N=p^rN',\ (p,N')=1)$. C'est donc une \uline{extension} \deleted{est} \uline{résoluble}, \inserted{mais mieux :} \deleted{\unsure{On en peut dire plus, sans étude approfondie. On obtient même la précision abélienne unitaire}} \deleted{Or} les $p$-\deleted{extensions} sous-groupes \inserted{$H$} de $\unsure{GP|G\mathbb P}(1,k)$ sont \uline{abéliens} et contenus dans \unsure{sous-groupe unipotent|sous-groupes unipotents} de dim $1$, \inserted{conjugués entre eux, et $H$ isom. à $(\mathbb Z/p\mathbb Z)^{\unsure{\nu|r}}$.} On trouve donc $G\subset\mathrm{Borel}$, et par suite $G$ produit $1/2$ direct d'un groupe cyclique premier à \unsure{2|?|p} car., par un $(\mathbb Z/p\mathbb Z)^r$. En fait, on voit aisément (\deleted{\unsure{si car $=$}} qu'on est soit dans le premier cas (cyclique premier à car, donc kummerien, qui n'est exclu), ou dans la \unsure{direction|deuxième}}

% slanted
\margin{\unsure{En tous les cas, on a une extension abélienne, soit kummerienne|En tous les cas, on a une extension abélienne, non kummerienne|En tous les cas, on a une extension résoluble, soit kummerienne} (d'allure \unsure{purement insép.|premier à $p$}), sauf \unsure{en type|que grpe} $(\mathbb Z/p^{\unsure{n|r}}\mathbb Z)^\nu$}

\deleted{\uline{Etude du cas} $(\deleted{\unsure{d}}N,2,2)$, $\{\,[G]=2N\,\}$. En car. $\ne2$, l'\uline{indice} \inserted{de} ramification $2$ \deleted{\unsure{de type modéré}} a ramification modérée, et faisant une extension kummerienne de degré $2$, on \deleted{\unsure{trou}} trouve un $X'/Y'$ qui n'a plus que $2$ pts de ramification \deleted{\unsure{dans}} $Y'$, donc $X'/Y'$ est justiciable de l'étude précédente mais \deleted{\unsure{il}} faut vérifier \deleted{\unsure{si tel est si}} \unsure{que} $X'$ de genre $0$ (facile)}

\diagram{\begin{tikzcd}[ampersand replacement=\&,column sep=small,row sep=small]
X \& X'\arrow[l,dash]\arrow[d,dash]\\
Y \& Y'\arrow[l,dash]
\arrow[from=1-1,to=2-1,dash]
\end{tikzcd}}
% archive sheet number: 67
